% BOOK-FILE-ID: f-c0e3f5877230
\subsection{解法のつながり}\label{節：解法のつながり}
既習の操作を別の対象へ使う7題である.
問1・2は置く量,問3・4は添字の扱い,問5〜7は周期・近似・和を扱う.
本文と関係が深いのは問1・2・4・5,\,\chapterref*{章：解法の背景と広がり}と関係が深いのは問3・6・7である.
読んだ説明に対応する題を選べる.

\subsubsection{問題}
\begin{enumerate}[label=\textbf{問\arabic*.},format=\bookitemlink{practice-connections}{problem}{answer}]
  \item $a_1=3$,
    $a_{n+1}=\frac{n+2}{n}a_n+n(n+1)(n+2)$ $(n\ge1)$とする.
    $b_n=\frac{a_n}{n(n+1)}$が満たす漸化式を求め,\bookterm{general}{一般項}を求めよ.
    また,元の漸化式を満たす1つの\bookterm{sequence}{数列}$v_n$を見つけ,
    $a_n-v_n$を考える別解を説明せよ.
  \item $a_1=1$, $a_2=3$,
    $a_{n+2}=\frac{a_{n+1}^2+1}{a_n}$ $(n\ge1)$とする.
    各項が正であることを示し,
    $I_n=\frac{a_n^2+a_{n+1}^2+1}{a_na_{n+1}}$を使って\bookterm{general}{一般項}を求めよ.
    初めの二項を正の値$x,y$に変えると,\bookterm{invariant}{不変量}の値はどう変わるか.
  \item $n,m\ge0$とする.\,$A_{n,0}=A_{0,m}=1$,
    $A_{n,m}=A_{n-1,m}+A_{n,m-1}$ $(n,m\ge1)$から,
    $A_{n,m}=\binom{n+m}{m}$を$n+m$に関する帰納法で証明せよ.
    また右と上に進む格子経路の数として説明せよ.
  \item $a_1=2$とし,
    \[
      a_{n+1}=\begin{cases}
        -a_n&(n\text{が奇数}),\\
        2a_n+1&(n\text{が偶数})
      \end{cases}
    \]
    とする.奇数番目と偶数番目の\bookterm{general}{一般項}を求めよ.
    係数が周期的であることと,\bookterm{sequence}{数列}自身が周期的であることは同じか.
  \item $a_1=1$, $a_2=2$,
    $a_{n+2}=\frac{1+a_{n+1}}{a_n}$ $(n\ge1)$とする.
    周期を使って\bookterm{general}{一般項}を定め,最小周期を求めよ.
    正の初期値を変えて定\bookterm{sequence}{数列}となる場合も求めよ.
  \item $x_1=3$, $y_1=2$,
    $x_{k+1}=3x_k+4y_k$, $y_{k+1}=2x_k+3y_k$ $(k\ge1)$とする.
    $x_k^2-2y_k^2=1$を証明せよ.
    $r_k=\frac{x_k}{y_k}$について,\,$r_k-\sqrt2$を$y_k,r_k$で表して\bookterm{limit}{極限}を求め,
    $\frac12\left(r_k+\frac{2}{r_k}\right)=r_{2k}$も示せ.
  \item $j^3=6\binom j3+6\binom j2+\binom j1$を確かめよ.
    二項係数の\bookterm{difference}{階差}を使って,\,$\sum_{j=0}^{N-1}j^3$ $(N\ge0)$を求めよ.
    単項式の係数比較で求める方法と,何が違うか述べよ.
\end{enumerate}

\subsubsection{方針}
\begin{enumerate}[label=\textbf{問\arabic*.},format=\bookitemlink{practice-connections}{hint}{problem}]
  \item 割る因子の隣り合う比を確認する.残った\bookterm{difference}{階差}は$n$となる.
    両辺を割った後の式の1つの解から,引く量$v_n$を逆に作る.
  \item 正値性を先に帰納的に確かめる.\,$xz=y^2+1$から$I_{n+1}=I_n$を示し,
    初期値で係数を決める.
  \item 内側の二つの位置では添字の和が一つ小さい.
    境界を先に確かめ,二項係数の加法公式を使う.
  \item 2回分の漸化式をまとめる.\,$x_k=a_{2k-1}$は\bookterm{constant}{定数係数}の一次\bookterm{recurrence}{漸化式}になる.
  \item 最初の七項を計算し,直前の二項の組が戻ることを確かめる.
    \bookterm{initial}{初期条件}が変わる場合は文字で計算する.
  \item 保存される式は,漸化式を代入する前後で展開して比べる.
    誤差は$x_k^2-2y_k^2$の因数分解,添字の倍増は二次式の平方から得る.
  \item 連続整数の積へ展開して恒等式を確かめる.
    各二項係数の和は,上の添字$N$で下の添字を一つ増やしたものになる.
\end{enumerate}

\subsubsection{解答}\label{節：解法のつながり：解答}\label{節：解法のつながり：答え}
\begin{enumerate}[label=\textbf{問\arabic*.},format=\bookitemlink{practice-connections}{answer}{problem}]
  \item \textbf{チャートの経路：}数列の項の係数から指定の式で割り,階差型にする方法と一つの解を引く方法を比べる.\par\smallskip
    \textbf{結論}\par\nopagebreak[4]
$b_{n+1}-b_n=n$, $b_1=\frac32$より,
    \[
      a_n=\frac{n(n+1)(n^2-n+3)}2.
    \]
    $v_n=\frac{n^2(n-1)(n+1)}2$を選ぶと,
    $a_n-v_n=\frac32n(n+1)$となる.
\par\smallskip\noindent\textbf{計算・証明}\par\nopagebreak[4]
\textbf{両辺を割って階差型にする.}
    $h_n=n(n+1)$は$\frac{h_{n+1}}{h_n}=\frac{n+2}{n}$を満たす.
    両辺を$(n+1)(n+2)$で割ると$b_{n+1}=b_n+n$であり,
    $b_n=\frac32+\sum_{j=1}^{n-1}j=\frac{n^2-n+3}{2}$となる.
    \par\smallskip\textbf{引く方法との比較.}
    両辺を割った後の式では,\,$w_n=\frac{n(n-1)}{2}$も$w_{n+1}-w_n=n$を満たす.
    $v_n=h_nw_n$は元の漸化式を満たすため,
    $d_n=a_n-v_n$について$d_{n+1}=\frac{n+2}{n}d_n$となる.
    $d_1=3$から$d_n=\frac32n(n+1)$を得る.
    この場合は割る方法から引く量も発見できた.
    \sectionref*{節：変数係数の一次型}へ戻って比べよう.

  \item \textbf{チャートの経路：}項どうしの積から指定の不変量を確かめ,正値を用いて隣接3項型へ直す.\par\smallskip
    \textbf{結論}\par\nopagebreak[4]
$I_n=\frac{11}{3}$なので$a_{n+2}=\frac{11}{3}a_{n+1}-a_n$となる.
    $\alpha=\frac{11+\sqrt{85}}{6}$, $\beta=\frac{11-\sqrt{85}}{6}$とすると,
    \[
      a_n=\frac{(3-\beta)\alpha^{n-1}+(\alpha-3)\beta^{n-1}}{\alpha-\beta}.
    \]
    正の初期値$x,y$では$I_n=\frac{x^2+y^2+1}{xy}$となる.
\par\smallskip\noindent\textbf{計算・証明}\par\nopagebreak[4]
\textbf{正値と保存の確認.}
    正の二項から次も正になる.\,$x=a_n,y=a_{n+1},z=a_{n+2}$とすると,\,$xz=y^2+1$である.
    これを不変量の分子に使えば,
    \[
      I_n=\frac{x^2+y^2+1}{xy}
      =\frac{x^2+xz}{xy}=\frac{x+z}{y}.
    \]
    次の位置でも
    \[
      I_{n+1}=\frac{y^2+z^2+1}{yz}
      =\frac{xz+z^2}{yz}=\frac{x+z}{y}
    \]
    となるので,\,$I_{n+1}=I_n$である.
    \par\smallskip\textbf{定数係数の漸化式へ戻す.}
    $I_1=\frac{1+9+1}{3}=\frac{11}{3}$なので$x+z=\frac{11}{3}y$となる.
    特性方程式の2つの解で$a_n=C\alpha^{n-1}+D\beta^{n-1}$と置き,
    $C+D=1$, $C\alpha+D\beta=3$を解くと解答の係数になる.
    本文の例とは初期値が異なり,不変量の値も$3$とは異なる.
    $I_n$の値は初期値によって変わるが,初期値を一組選べば,その\bookterm{sequence}{数列}の中では一定である.
    \sectionref*{節：変わらない量を使う}へ.

  \item \textbf{チャートの経路：}式全体の二つの添字に注目し,添字の和による帰納法と格子経路の計数を使う.\par\smallskip
    \textbf{境界から内側へ.}
    $n=0$または$m=0$では式は$1$になる.
    $n,m\ge1$の場合,\,$n+m$より小さい添字の和で式が成立すると仮定し,
    \BookMathLayout{\[
        A_{n,m}=\binom{n+m-1}{m}+\binom{n+m-1}{m-1}
        =\binom{n+m}{m}
      \]}{\[
        \begin{aligned}
          A_{n,m}&=\binom{n+m-1}{m}+\binom{n+m-1}{m-1}\\
          &=\binom{n+m}{m}
        \end{aligned}
      \]}
    を得る.格子経路では最後の一歩が右か上かで分けた和であり,
    全歩の並べ方を直接数えても同じ二項係数になる.
    \sectionref*{節：二方向の階差}へ.

  \item \textbf{チャートの経路：}係数の周期から2回分の式をまとめ,奇数番目を特性方程式型にする.\par\smallskip
    \textbf{結論}\par\nopagebreak[4]
$k\ge1$について,
    \BookMathLayout{\[
        a_{2k-1}=\frac{1+5(-2)^{k-1}}3,\qquad
        a_{2k}=-\frac{1+5(-2)^{k-1}}3.
      \]}{\[
        \begin{aligned}
          a_{2k-1}&=\frac{1+5(-2)^{k-1}}3,\\
          a_{2k}&=-\frac{1+5(-2)^{k-1}}3.
        \end{aligned}
      \]}
    絶対値が非有界になるので,\bookterm{sequence}{数列}自身は周期的ではない.
\par\smallskip\noindent\textbf{計算・証明}\par\nopagebreak[4]
\textbf{一周期分をまとめる.}
    $x_k=a_{2k-1}$とすると,\,$a_{2k}=-x_k$から
    $x_{k+1}=-2x_k+1$となる.
    \bookterm{fixed}{不動点}$\frac{1}{3}$を引けば$x_{k+1}-\frac{1}{3}=-2\left(x_k-\frac{1}{3}\right)$であり,
    $x_1=2$から解答の式が得られる.
    係数の周期が$2$でも,一周期後の係数$-2$によって項の絶対値は大きくなる.
    \sectionref*{節：一周期分をまとめる}へ.

  \item \textbf{チャートの経路：}非線形な3項間の式で正値を確かめ,2項の組の繰り返しから周期を求める.\par\smallskip
    \textbf{結論}\par\nopagebreak[4]
各項は$1,2,3,2,1$を繰り返す.
    $n-1$を$5$で割った余りが$0,1,2,3,4$のとき,
    $a_n$はそれぞれ$1,2,3,2,1$である.最小周期は$5$.
    正の定\bookterm{sequence}{数列}は$a_1=a_2=\frac{1+\sqrt5}{2}$の場合である.
\par\smallskip\noindent\textbf{計算・証明}\par\nopagebreak[4]
\textbf{二項の組が戻る.}
    $a_3=3,a_4=2,a_5=1,a_6=1,a_7=2$なので,
    $(a_6,a_7)=(a_1,a_2)$である.連続する2項から以降の項はすべて決まるので,全体が繰り返す.
    周期$5$の最小周期は$1$または$5$であり,定\bookterm{sequence}{数列}ではないので$5$である.
    定数$t>0$では$t^2=t+1$が必要十分である.
    \sectionref*{節：非線形の周期性}へ.

  \item \textbf{チャートの経路：}連立の平方差を不変量として確かめ,積と比の関係から誤差と添字の倍増を調べる.\par\smallskip
    \textbf{結論}\par\nopagebreak[4]
$x_k^2-2y_k^2=1$が保たれ,
    \[
      r_k-\sqrt2=\frac1{y_k^2(r_k+\sqrt2)},\qquad
      \lim_{k\to\infty}r_k=\sqrt2.
    \]
    また$r_{2k}=\frac12\left(r_k+\frac{2}{r_k}\right)$である.
\par\smallskip\noindent\textbf{計算・証明}\par\nopagebreak[4]
\textbf{保存・近似・添字の倍増.}
    漸化式を代入して展開すれば
    \[
      (3x+4y)^2-2(2x+3y)^2=x^2-2y^2
    \]
    となり,初期値での値$1$が保存される.
    正値性と$y_{k+1}>3y_k$から$y_k\to\infty$であり,
    $(r_k-\sqrt2)(r_k+\sqrt2)=\frac{1}{y_k^2}$から誤差と\bookterm{limit}{極限}が得られる.
    また漸化式と初期値から$x_k+y_k\sqrt2=(3+2\sqrt2)^k$である.
    平方して$x_{2k}=x_k^2+2y_k^2$, $y_{2k}=2x_ky_k$を得れば,
    比は\bookterm{newton}{ニュートン法}の式になる.
    \sectionref*{節：整数から作る近似分数}へ.

  \item \textbf{チャートの経路：}和を求めるため多項式を二項係数で表し,隣接する二項係数の差を使う.\par\smallskip
    \textbf{結論}\par\nopagebreak[4]
\BookMathLayout{\[
        \sum_{j=0}^{N-1}j^3
        =6\binom N4+6\binom N3+\binom N2
        =\frac{N^2(N-1)^2}{4}.
      \]}{\[
        \begin{aligned}
          \sum_{j=0}^{N-1}j^3&=6\binom N4+6\binom N3+\binom N2\\
          &=\frac{N^2(N-1)^2}{4}.
        \end{aligned}
      \]}
    階差を取りやすい部品へ表し直してから,部品ごとに和を取った.
\par\smallskip\noindent\textbf{計算・証明}\par\nopagebreak[4]
\textbf{操作に合う表示.}
    $j(j-1)(j-2)+3j(j-1)+j=j^3$だから,提示された恒等式が成り立つ.
    階差の和の公式をそれぞれに使うと,
    $6\binom N4+6\binom N3+\binom N2$となる.
    $N\ge4$では階乗表示を整理して解答の式を得る.
    $N=0,1,2,3$でも,範囲外の二項係数を0とする規約で直接確かめられる.
    係数比較で差を逆に戻す方法に対し,ここでは階差を取ると添字が1つ下がる部品を使った.
    \sectionref*{節：二項係数で多項式を表す}へ.
\end{enumerate}

\subsubsection{解説}\label{節：解法のつながり：解説}
解法1は解答で用いた方法である.以下では,その選択の理由や別解との違いを説明する.
\begin{enumerate}[label=\textbf{問\arabic*.},format=\bookitemlink{practice-connections}{solution}{problem}]
  \item \textbf{【解法の選択】}\par\nopagebreak[4]
両辺を割る方法では,隣接する比が係数$\frac{n+2}{n}$になる$n(n+1)$を選ぶ.
一つの解を引く方法では,元の漸化式全体を満たす$v_n$を探す.
\par\smallskip\noindent\textbf{【比較】}\par\nopagebreak[4]
割った後の階差型で一つの解を作り,それに$n(n+1)$を掛ければ,引く数列$v_n$を作れる.
二つの方法は独立した技巧ではなく,同じ係数を処理した後,付加項を和で扱うか一つの解として先に引くかの違いである.
設問で要求された両方法の計算は,解答に示した.
\end{enumerate}

\begin{enumerate}[start=6,label=\textbf{問\arabic*.},format=\bookitemlink{practice-connections}{solution}{problem}]
  \item \textbf{【解法の選択】}\par\nopagebreak[4]
平方差が一定であることは代入で確認できる.これを因数分解すれば近似値の誤差を評価できる.
添字が倍になる式は,元の二つの数列を使って考えよう.
\par\smallskip\noindent\textbf{【もう一つの導き方】}\par\nopagebreak[4]
$z_k=x_k+\sqrt2y_k$とおくと,\,$z_{k+1}=(3+2\sqrt2)z_k$である.
初項から$z_k=(3+2\sqrt2)^k$なので,\,$z_{2k}=z_k^2$が成り立つ.
共役な式も合わせて実数部分と$\sqrt2$の係数を比べれば,
\begin{align*}
  x_{2k}&=x_k^2+2y_k^2,\\
  y_{2k}&=2x_ky_k.
\end{align*}
したがって
\begin{align*}
  r_{2k}&=\frac{x_k^2+2y_k^2}{2x_ky_k}\\
        &=\frac12\left(r_k+\frac2{r_k}\right).
\end{align*}
\par\smallskip\noindent\textbf{【比較】}\par\nopagebreak[4]
元の整数列で添字を倍にすることは,その比にニュートン法を1回使うことに対応する.
整数性や不変量を示すには元の連立式が便利であり,近似値の改善を読むには比の式が便利である.

  \item \textbf{【解法の選択】}\par\nopagebreak[4]
二項係数では隣接する差を取ると下の添字が一つ下がる.
そのため,\,$j^3$を二項係数で表せば,和を取る操作を逆向きに行える.
\par\smallskip\noindent\textbf{【別の方法との比較】}\par\nopagebreak[4]
通常の係数比較なら,和を$AN^4+BN^3+CN^2+DN$と置き,隣接する差が$N^3$になるよう四つの係数を決める.
二項係数を使う方法では,差の取り方が既知なので,和を取るたびに係数を解き直す必要がない.
ただし,最初に多項式を二項係数で表す計算が必要である.どちらも同じ和を求めるが,計算しやすい式の選び方が異なる.
\end{enumerate}
